\section*{CHƯƠNG 1. MỞ ĐẦU}
\addcontentsline{toc}{section}{\numberline{}CHƯƠNG 1. MỞ ĐẦU}
\setcounter{section}{1}
\setcounter{subsection}{0}
\setcounter{figure}{0}
\setcounter{table}{0}

\subsection{Bối cảnh và tính cấp thiết của đề tài}
Trong ngành dịch vụ và bán lẻ hiện đại, việc thấu hiểu hành vi và cảm xúc của khách hàng đã trở thành nền tảng cho sự tăng trưởng kinh doanh bền vững. Khi thị trường chuyển dịch từ định hướng sản phẩm sang định hướng lấy khách hàng làm trung tâm, các tổ chức liên tục tìm kiếm những phương pháp đổi mới sáng tạo để đo lường, phân tích và cải thiện Trải nghiệm Khách hàng (CX). Các phương pháp đánh giá mức độ hài lòng truyền thống — như khảo sát thủ công, bảng câu hỏi và phỏng vấn sau dịch vụ — luôn tồn tại những hạn chế cố hữu. Những phương pháp mang tính hồi cứu này thường tốn thời gian, gây phiền toái và dễ bị ảnh hưởng bởi thiên kiến nhận thức, do khách hàng có thể không nhớ chính xác hoặc không báo cáo trung thực trạng thái cảm xúc của họ. Hệ quả là, doanh nghiệp gặp khó khăn trong việc nắm bắt các phản ứng cảm xúc nguyên bản, tức thời và vô thức của khách hàng trong quá trình họ tương tác với sản phẩm, quảng cáo hoặc môi trường dịch vụ.

Để giải quyết những thách thức này, sự tiến bộ vượt bậc của Trí tuệ Nhân tạo (AI), đặc biệt là Thị giác Máy tính và Học sâu, đã mở ra những mô hình mới cho việc phân tích khách hàng một cách tự động và không xâm lấn. Biểu cảm khuôn mặt, với tư cách là phương tiện giao tiếp phi ngôn ngữ chính của con người, mang thông tin phong phú và trực tiếp về trạng thái nhận thức và cảm xúc của một cá nhân. Các hệ thống Nhận dạng Biểu cảm Khuôn mặt (FER) tận dụng camera kỹ thuật số và mạng nơ-ron sâu để tự động phát hiện khuôn mặt, trích xuất các đặc trưng chi tiết và phân loại các trạng thái cảm xúc theo thời gian thực. Bằng cách triển khai các hệ thống dựa trên FER trong môi trường bán lẻ hoặc tại các điểm chạm kỹ thuật số, doanh nghiệp có thể quan sát động lực cảm xúc của khách hàng một cách tự nhiên, cho phép họ đánh giá mức độ tương tác với sản phẩm, đo lường hiệu quả quảng cáo và đánh giá chất lượng dịch vụ tổng thể một cách định lượng và khách quan.

Do đó, việc phát triển một hệ thống phân tích khách hàng mạnh mẽ và theo thời gian thực dựa trên biểu cảm khuôn mặt không chỉ có ý nghĩa khoa học lớn mà còn mang lại giá trị thực tiễn vô cùng to lớn. Nó cung cấp cho các doanh nghiệp trí tuệ cảm xúc dựa trên dữ liệu để tối ưu hóa các chiến lược tiếp thị, cá nhân hóa dịch vụ chăm sóc khách hàng và cuối cùng là nâng cao mức độ hài lòng của khách hàng.

\subsection{Đặt vấn đề}
Mặc dù nhận dạng biểu cảm khuôn mặt có tiềm năng lớn trong phân tích khách hàng, nhiều thách thức kỹ thuật và thực tiễn vẫn chưa được giải quyết triệt để khi triển khai trong thế giới thực:
\begin{itemize}
    \item \textbf{Sự biến động của môi trường:} Trong môi trường bán lẻ thực tế, hình ảnh khuôn mặt thu được từ camera giám sát tiêu chuẩn hoặc ki-ốt thường bị giảm chất lượng bởi nhiều yếu tố phức tạp, bao gồm thay đổi ánh sáng môi trường, các góc chụp khác nhau (biến động tư thế), nhòe do chuyển động và bị che khuất một phần (như kính mắt, khẩu trang hoặc tay). Việc phát triển một mô hình học sâu đủ mạnh mẽ để tổng quát hóa tốt dưới các điều kiện không kiểm soát này là một thách thức kỹ thuật lớn.
    \item \textbf{Sự khác biệt cá nhân và văn hóa:} Cường độ biểu cảm và cấu trúc khuôn mặt thay đổi rất lớn giữa các nhóm tuổi, giới tính và chủng tộc khác nhau. Một mô hình được huấn luyện trên các tập dữ liệu công khai chung có thể bị suy giảm hiệu năng khi triển khai trên các nhóm nhân khẩu học bản địa cụ thể, đòi hỏi các kỹ thuật trích xuất đặc trưng có khả năng thích ứng cao.
    \item \textbf{Yêu cầu xử lý theo thời gian thực:} Đối với các quầy dịch vụ khách hàng hoặc bảng hiệu kỹ thuật số tương tác, phản hồi cảm xúc cần được xử lý ngay lập tức. Việc cân bằng giữa độ chính xác của mô hình và hiệu quả tính toán để đạt được tốc độ suy luận thời gian thực trên các thiết bị tính toán biên hoặc phần cứng có ngân sách hạn chế là vô cùng quan trọng.
    \item \textbf{Tạo ra các thông tin chi tiết hữu ích:} Các nhãn cảm xúc thô (ví dụ: vui vẻ, trung tính, ngạc nhiên) đứng độc lập là chưa đủ cho việc ra quyết định kinh doanh. Xây dựng một khung phân tích giúp tổng hợp các biểu cảm khuôn mặt theo thời gian thành các chỉ số kinh doanh cấp cao, như Điểm tương tác của khách hàng (CES) hoặc Chỉ số hài lòng của khách hàng (CSI), là vô cùng cần thiết.
\end{itemize}

\subsection{Mục tiêu nghiên cứu}
Đồ án tốt nghiệp này hướng tới việc thiết kế, triển khai và đánh giá một hệ thống phân tích khách hàng tự động dựa trên biểu cảm khuôn mặt. Để đạt được điều này, nghiên cứu tập trung vào các mục tiêu cụ thể sau:
\begin{enumerate}
    \item Thiết kế một quy trình phần mềm hoàn chỉnh và tối ưu hóa để phân tích khuôn mặt theo thời gian thực, bao gồm phát hiện khuôn mặt, xác định vị trí các điểm mốc khuôn mặt, căn chỉnh và phân loại biểu cảm khuôn mặt.
    \item Phát triển và huấn luyện một mô hình học sâu mạnh mẽ (ví dụ: sử dụng Mạng nơ-ron tích chập - CNN hoặc các mô hình Vision Transformer gọn nhẹ) có khả năng phân loại biểu cảm khuôn mặt thành các danh mục cảm xúc cơ bản trong các điều kiện môi trường không bị ràng buộc.
    \item Xây dựng một hệ thống phân tích cảm xúc giúp chuyển đổi các chuỗi biểu cảm theo thời gian thực thành các chỉ số định lượng về sự hài lòng và mức độ tương tác của khách hàng.
    \item Thực hiện các thực nghiệm toàn diện trên các tập dữ liệu công khai và các kịch bản thực tế để kiểm chứng độ chính xác, độ trễ và tính khả thi của hệ thống trong các ứng dụng kinh doanh thực tiễn.
\end{enumerate}

\subsection{Phạm vi nghiên cứu}
Phạm vi của nghiên cứu này được xác định như sau:
\begin{itemize}
    \item \textbf{Mô hình cảm xúc:} Hệ thống sẽ phân loại các biểu cảm thành bảy trạng thái cảm xúc cơ bản dựa trên lý thuyết cảm xúc của Ekman: Vui vẻ, Buồn bã, Tức giận, Ngạc nhiên, Sợ hãi, Ghê tởm và Trung tính.
    \item \textbf{Kịch bản ứng dụng:} Hệ thống được thiết kế cho môi trường kinh doanh trong nhà, chẳng hạn như quầy thanh toán của cửa hàng bán lẻ, màn hình quảng cáo kỹ thuật số hoặc quầy thử nghiệm sản phẩm, sử dụng camera tiêu chuẩn hướng thẳng diện hoặc hơi nghiêng phía trên.
    \item \textbf{Đối tượng phân tích:} Phân tích được thực hiện ở cấp độ tổng hợp ẩn danh để tôn trọng quyền riêng tư của khách hàng, tập trung vào các xu hướng cảm xúc và chỉ số thống kê thay vì theo dõi danh tính từng cá nhân.
\end{itemize}

\subsection{Bố cục đồ án}
Phần còn lại của đồ án được cấu trúc như sau:
\begin{itemize}
    \item \textbf{Chương 2: Cơ sở lý thuyết và Tổng quan tài liệu} cung cấp một cái nhìn sâu sắc về công nghệ phân tích khuôn mặt, các kiến trúc học sâu cho bài toán nhận dạng biểu cảm khuôn mặt (FER), và các ứng dụng thực tế của AI cảm xúc trong kinh doanh.
    \item \textbf{Chương 3: Phương pháp nghiên cứu và Thiết kế hệ thống} chi tiết hóa kiến trúc hệ thống đề xuất, thiết kế mô hình học sâu, quy trình tiền xử lý dữ liệu và khung phân tích cảm xúc khách hàng.
    \item \textbf{Chương 4: Kết quả thực nghiệm và Thảo luận} trình bày về thiết lập thực nghiệm, quy trình huấn luyện, các chỉ số đánh giá, phân tích hiệu năng và kết quả thử nghiệm triển khai thực tế.
    \item \textbf{Chương 5: Kết luận và Hướng phát triển} tóm tắt các thành tựu chính của đồ án, thảo luận về các hạn chế và đề xuất các hướng nghiên cứu trong tương lai.
\end{itemize}